\section{Related Work}
\label{sec:related_work}

The intersection of deep learning compilers and energy optimization is an active area of research. Early tensor compilers, such as TVM and its automated scheduling framework Ansor \cite{zheng2020ansor}, pioneered the use of machine learning to explore the vast configuration space of tensor programs. However, Ansor and similar frameworks primarily focus on execution latency.

Recently, several studies have attempted to integrate energy as an optimization objective. Schoonhoven et al. \cite{schoonhoven2022going} explored model-steered autotuning for GPU energy efficiency, but their approach required a custom compiler backend rather than integrating with existing, widely-adopted DSLs like Triton. Cheema and Khan \cite{cheema2023gpuautotuning} developed an autotuning framework for OpenCL kernels, demonstrating that multi-objective search can yield significant power savings. However, OpenCL is rarely used in modern LLM inference stacks.

More closely related to our work, Zhang et al. \cite{zhang2024automating} introduced a fast search-based compilation approach for energy-efficient GPU kernel generation. While effective, their methodology requires training a surrogate cost model offline. In contrast, \textbf{GreenTune} operates directly on the device, providing an online, energy-aware extension to Triton's native autotuner without the need for pre-trained models.

Finally, emerging studies on LLM-based kernel generation, such as KernelPro \cite{kernelpro2026} and FlipFlop \cite{flipflop2026}, treat energy as a secondary metric (a "tie-breaker" subordinate to speedup) or rely on static analysis. Our work differs by promoting energy to a first-class objective within the runtime autotuner of the most ubiquitous tensor compiler in modern AI systems.
